\documentclass[11pt]{article}

\usepackage[a4paper,margin=0.82in]{geometry}
\usepackage[T1]{fontenc}
\usepackage[utf8]{inputenc}
\usepackage{lmodern}
\usepackage{amsmath,amssymb,amsthm,mathtools}
\usepackage{booktabs,enumitem,microtype,tabularx,xcolor}
\usepackage[numbers,sort&compress]{natbib}
\usepackage[colorlinks=true,linkcolor=blue!65!black,
  citecolor=blue!65!black,urlcolor=blue!70!black]{hyperref}
\usepackage[nameinlink,noabbrev]{cleveref}
\hypersetup{
  pdftitle={Resolved Principal Mass Before Width Grouping},
  pdfauthor={Liang Wang},
  pdfsubject={TPC-104 exact endpoint and inverse-length disintegration},
  pdfkeywords={principal mass, resonance width, inverse length, fixed shift, opened atoms}
}

\newtheorem{theorem}{Theorem}[section]
\newtheorem{proposition}[theorem]{Proposition}
\newtheorem{corollary}[theorem]{Corollary}
\newtheorem{lemma}[theorem]{Lemma}
\theoremstyle{definition}
\newtheorem{definition}[theorem]{Definition}
\newtheorem{example}[theorem]{Example}
\theoremstyle{remark}
\newtheorem{remark}[theorem]{Remark}

\newcommand{\Z}{\mathbb Z}
\newcommand{\one}{\mathbf 1}
\newcommand{\Lzero}{\mathrm{L0}}
\newcommand{\Lone}{\mathrm{L1}}
\newcommand{\Ltwo}{\mathrm{L2}}
\newcommand{\PX}{\mathfrak P_X}
\newcommand{\EX}{\mathfrak E_X}
\newcommand{\LX}{\mathfrak L_X}

\title{\textbf{Resolved Principal Mass Before Width Grouping:}\\
Exact Endpoint Return and an Inverse-Length Equivalence}
\author{Liang Wang\\
School of Artificial Intelligence and Automation\\
Huazhong University of Science and Technology, Wuhan 430074, P.R. China\\
\texttt{wangliang.f@gmail.com}}
\date{July 2026}

\begin{document}
\maketitle

\begin{abstract}
The positive multiplicative-resonance route contains the unresolved
principal mass
\[
 \mathfrak P_X
 =
 \sum_\beta
 \frac{2H_{q_X}(N_\beta)}{q_X-1}
 \sum_{\substack{z\in I_\beta^\ast\\q_X\nmid u_{\beta,z}}}
 w_{\beta,z},
 \qquad
 H_q(N)=
 \min\!\left\{\frac{q-1}{2},\left\lfloor\frac qN\right\rfloor\right\}.
\]
We open this quantity before grouping equal widths or multipliers
and obtain an exact endpoint/long-fiber disintegration.

For \(N=1,2\), one has
\[
 \frac{2H_q(N)}{q-1}=1,\qquad
 g_q(H_q(N))=0.
\]
Thus every remaining invertible atom on these fibers is purely
principal and enters with its full physical weight.  For
\(3\le N\le q\), the elementary but sharp uniform comparison
\[
 \frac1N
 \le
 \frac{2\lfloor q/N\rfloor}{q-1}
 \le
 \frac{2q}{(q-1)N}
\]
holds.  If
\[
 \mathfrak E_X
 =
 \sum_{\substack{\beta:N_\beta\le2\\z\in I_\beta^\ast\\
                  q_X\nmid u_{\beta,z}}}w_{\beta,z},
 \qquad
 \mathfrak L_X
 =
 \sum_{\substack{\beta:3\le N_\beta\le q_X}}
 \frac1{N_\beta}
 \sum_{\substack{z\in I_\beta^\ast\\q_X\nmid u_{\beta,z}}}
 w_{\beta,z},
\]
then
\[
 \boxed{
 \mathfrak E_X+\mathfrak L_X
 \le\mathfrak P_X
 \le
 \mathfrak E_X+
 \frac{2q_X}{q_X-1}\mathfrak L_X.}
\]
Consequently the principal gate is equivalent, up to absolute
constants and in the same physical normalization, to controlling
the actual endpoint mass and inverse-length mass separately.

This is an exact \(\Lone\) reduction, not a proof of either bound.
We also give sharp positive countermodels showing that the
\(1/N_\beta\) factor alone cannot control a polynomial number of
branches, and that a long-fiber estimate cannot return the
\(N=1,2\) mass.  No \(\Ltwo\) fixed-shift estimate or parity
breakthrough is claimed.
\end{abstract}

\noindent\textbf{Keywords:}
principal resonance; resolved length; endpoint fiber; inverse-length
mass; positive majorant; fixed shift.

\medskip
\noindent\textbf{Literal-frame convention.}
Every sum uses the actual \(q_X\nmid u\) opened atoms, their physical
weights, one prescribed \(h_0\ne0\), and one inherited global
normalization.  No ambient branch count replaces an opened mass,
and nothing is specialized to \(h_0=2\).

\section{The unresolved principal term}
\label{sec:introduction}

TPC-99 represents nonconstant resonance by the incidence window
\[
 E_H=\{\pm1,\ldots,\pm H\}\subset\mathbb F_q^\times.
\]
TPC-100 then opens each frequency-free branch
\(\beta=(\theta,c,\kappa)\) into atoms
\[
 p=(\beta,z),\qquad
 w_{\beta,z}\ge0,\qquad
 w_\beta=\sum_{z\in I_\beta^\ast}w_{\beta,z},
\]
without changing its resolved length
\[
 N_\beta=\#I_\beta^\ast
\]
or its resonance width \(H_q(N_\beta)\)
\citep{WangTPC99,WangTPC100}.  After separately returning the
\(q_X\mid u_{\beta,z}\) submass, the remaining principal ledger is
\begin{equation}
 \PX
 =
 \sum_{\substack{\beta\\1\le N_\beta\le q_X}}
 \rho_{q_X}(N_\beta)
 \sum_{\substack{z\in I_\beta^\ast\\q_X\nmid u_{\beta,z}}}
 w_{\beta,z},
\qquad
 \rho_q(N):=\frac{2H_q(N)}{q-1}.
\label{eq:P}
\end{equation}
The frequency-filter factor \(B_X=\|b_X\|_1\ll1\) is not included
in \(\PX\).

TPC-102 proposes opening \eqref{eq:P} before grouping by equal
width.  That order matters.  Grouping first hides the fact that
full-width resonance occurs at \(N=1,2\) and
that the remaining coefficient is exactly comparable to \(1/N\)
\citep{WangTPC101,WangTPC102}.

\section{The exact width coefficient}
\label{sec:coefficient}

Let \(q\) be an odd prime and \(1\le N\le q\).  Recall
\[
 H_q(N)=
 \min\!\left\{\frac{q-1}{2},\left\lfloor\frac qN\right\rfloor\right\}.
\]

\begin{lemma}[The \(N=1,2\) full-width endpoint lengths]
\label{lem:endpoints}
For \(N=1,2\),
\[
 H_q(N)=\frac{q-1}{2},
 \qquad
 \rho_q(N)=1.
\]
For \(3\le N\le q\),
\[
 H_q(N)=\left\lfloor\frac qN\right\rfloor.
\]
\end{lemma}

For \(q\ge5\), these are the only full-width lengths.  At the
degenerate prime \(q=3\), \(N=3\) also has
\(H_3(3)=(q-1)/2\); it is nevertheless included correctly in the
\(N\ge3\) inverse-length ledger below.

\begin{proof}
For \(N=1\), the minimum in \(H_q(N)\) is visibly
\((q-1)/2\).  Since \(q\) is odd,
\(\lfloor q/2\rfloor=(q-1)/2\), proving the same assertion for
\(N=2\).  If \(N\ge3\), then
\[
 \left\lfloor\frac qN\right\rfloor
 \le\left\lfloor\frac q3\right\rfloor
 \le\frac{q-1}{2}.
\]
\end{proof}

\begin{lemma}[Sharp inverse-length comparison]
\label{lem:comparison}
For every odd prime \(q\) and \(3\le N\le q\),
\begin{equation}
 \frac1N
 \le\rho_q(N)
 \le\frac{2q}{(q-1)N}
 \le\frac3N.
\label{eq:comparison}
\end{equation}
The first inequality also holds for \(N=1,2\).
\end{lemma}

\begin{proof}
Write \(k=\lfloor q/N\rfloor\) and \(q=kN+r\) with
\(0\le r<N\).  Then
\[
 2Nk-(q-1)=kN-r+1\ge N-r+1>0.
\]
Thus
\[
 \rho_q(N)=\frac{2k}{q-1}\ge\frac1N.
\]
The upper bound follows from \(k\le q/N\).  Finally
\(2q/(q-1)\le3\) for \(q\ge3\).  At \(N=1,2\), use
\cref{lem:endpoints}.
\end{proof}

\begin{remark}[Optimal constants]
The lower constant \(1\) cannot be replaced uniformly by a number
larger than \(1\): for \(N=(q+1)/2\),
\[
 N\rho_q(N)=\frac{q+1}{q-1}\longrightarrow1.
\]
The asymptotic upper constant \(2\) is also sharp, for example at
\(N=q\), where \(N\rho_q(N)=2q/(q-1)\).
\end{remark}

\section{Resolved disintegration}
\label{sec:disintegration}

For one branch define its actual remaining invertible mass
\[
 M_\beta^{\mathrm{inv}}
 =
 \sum_{\substack{z\in I_\beta^\ast\\q_X\nmid u_{\beta,z}}}
 w_{\beta,z}.
\]
This may be smaller than the full branch mass \(w_\beta\); no
replacement by \(N_\beta\) or an ambient count is made.

\begin{definition}[Endpoint and inverse-length ledgers]
\label{def:ledgers}
Set
\begin{align}
 \EX
 &=
 \sum_{\substack{\beta\\N_\beta=1,2}}
 M_\beta^{\mathrm{inv}},
\label{eq:E}\\
 \LX
 &=
 \sum_{\substack{\beta\\3\le N_\beta\le q_X}}
 \frac{M_\beta^{\mathrm{inv}}}{N_\beta}.
\label{eq:L}
\end{align}
We also write
\[
 \EX=\mathfrak E_{1,X}+\mathfrak E_{2,X}
\]
for the exact \(N=1\) and \(N=2\) subledgers.
\end{definition}

\begin{theorem}[Exact endpoint/long-fiber disintegration]
\label{thm:disintegration}
The actual principal mass satisfies the identity
\begin{equation}
 \PX
 =
 \mathfrak E_{1,X}+\mathfrak E_{2,X}
 +
 \sum_{\substack{\beta\\3\le N_\beta\le q_X}}
 \frac{2\lfloor q_X/N_\beta\rfloor}{q_X-1}
 M_\beta^{\mathrm{inv}}.
\label{eq:exact-disintegration}
\end{equation}
Consequently,
\begin{equation}
 \boxed{
 \EX+\LX
 \le\PX
 \le
 \EX+\frac{2q_X}{q_X-1}\LX
 \le\EX+3\LX.}
\label{eq:equivalence}
\end{equation}
\end{theorem}

\begin{proof}
Insert \cref{lem:endpoints} into \eqref{eq:P} for
\(N_\beta=1,2\), and insert its long-fiber formula for
\(N_\beta\ge3\).  This gives
\eqref{eq:exact-disintegration}.  Apply
\cref{lem:comparison} term by term to the nonnegative literal
weights.
\end{proof}

\begin{corollary}[Equivalent formulation of Gate P]
\label{cor:gate}
The following are equivalent:
\begin{align*}
 \PX&\le X^{o(1)}Q_X^2,\\
 \EX+\LX&\le X^{o(1)}Q_X^2,\\
 \EX&\le X^{o(1)}Q_X^2
 \quad\text{and}\quad
 \LX\le X^{o(1)}Q_X^2.
\end{align*}
All three statements have the same quantifier scope and physical
normalization.
\end{corollary}

\begin{proof}
Use \eqref{eq:equivalence}, nonnegativity, and the fact that fixed
constants are absorbed in \(X^{o(1)}\).
\end{proof}

\begin{remark}[What ``equivalent'' does not mean]
\Cref{cor:gate} is an exact reduction of one sufficient positive
gate.  It does not identify that gate with the original signed
low-window sum.  A failure of \(\EX\) or \(\LX\) may stop the
termwise positive route while signed frequency or branch
cancellation remains possible.
\end{remark}

\section{The endpoint mass is exactly principal}
\label{sec:endpoint}

TPC-101 defines
\[
 g_q(H)=
 \sqrt{\frac{2H(q-1-2H)}{q-1}}.
\]

\begin{proposition}[Exact \(N=1,2\) return]
\label{prop:endpoint-return}
For \(N_\beta=1,2\),
\[
 g_{q_X}(H_{q_X}(N_\beta))=0.
\]
Moreover, for the fixed positive frequency weight \(b_X\),
the endpoint contribution to the multiplicative incidence carrier
is exactly
\begin{equation}
 B_X\EX,\qquad B_X=\sum_{r\in\mathbb F_{q_X}^\times}b_X(r).
\label{eq:endpoint-carrier}
\end{equation}
\end{proposition}

\begin{proof}
By \cref{lem:endpoints}, \(H=(q_X-1)/2\), so the displayed formula
for \(g_q\) vanishes.  At this width,
\[
 E_H=\mathbb F_{q_X}^\times.
\]
Hence the incidence indicator is one for every pair
\((r,\omega)\), and its bilinear form is the product of the two
total masses.  The opened census mass at these two lengths is
\(\EX\).
\end{proof}

\begin{remark}[No centered theorem can pay the endpoint]
The vanishing of \(g_q(H)\) is not a saving on \(\EX\).  It says
that the incidence operator has only its principal contribution at
full width.  Every surviving endpoint atom therefore has to be
returned through \eqref{eq:endpoint-carrier}, with its full physical
weight.
\end{remark}

\section{Consequences of the atom cap}
\label{sec:atom-cap}

TPC-103 proves on the same actual support that
\[
 W_X:=\max_{\beta,z}w_{\beta,z}=X^{o(1)}.
\]
This is the literal pointwise result of TPC-103
\citep{WangTPC103}.
Let
\[
 \mathcal B_{\le2,X}^{\mathrm{inv}}
 =
 \{\beta:N_\beta\le2,\ M_\beta^{\mathrm{inv}}>0\},
\qquad
 \mathcal B_{\ge3,X}^{\mathrm{inv}}
 =
 \{\beta:3\le N_\beta\le q_X,\
 M_\beta^{\mathrm{inv}}>0\}.
\]

\begin{proposition}[Exact-count sufficient conditions]
\label{prop:count}
One has
\begin{align}
 \EX
 &\le2W_X
 |\mathcal B_{\le2,X}^{\mathrm{inv}}|,
\label{eq:E-count}\\
 \LX
 &\le W_X
 |\mathcal B_{\ge3,X}^{\mathrm{inv}}|.
\label{eq:L-count}
\end{align}
Thus a subpower atom cap plus an
\(X^{o(1)}Q_X^2\) bound for each actual branch census is sufficient
for Gate P.
\end{proposition}

\begin{proof}
A branch of length \(N_\beta\) contains at most \(N_\beta\) opened
atoms.  Hence
\[
 M_\beta^{\mathrm{inv}}\le N_\beta W_X.
\]
For \(N_\beta\le2\), sum this inequality directly.  For
\(N_\beta\ge3\), divide by \(N_\beta\) before summing.
\end{proof}

\begin{remark}[Why this is not an ambient-count replacement]
\Cref{prop:count} is only a sufficient corollary after the exact
weighted identities have been proved.  It does not authorize
replacing \(M_\beta^{\mathrm{inv}}\) by \(N_\beta\) in
\eqref{eq:exact-disintegration}.  A future proof should use the
literal weights whenever they improve the crude branch census.
\end{remark}

\section{Sharp positive obstructions}
\label{sec:obstructions}

\begin{proposition}[Inverse length does not eliminate branch count]
\label{prop:branch-count}
For every \(M\ge1\) and \(3\le N\le q\), there is an abstract
nonnegative resolved census with \(M\) branches of length \(N\),
unit atom cap, and
\[
 \LX=M,\qquad
 \PX=M\,N\rho_q(N)\asymp M.
\]
Thus the \(1/N\) factor alone gives no bound independent of the
number of active branches.
\end{proposition}

\begin{proof}
Give every branch exactly \(N\) atoms of weight one.  Then
\(M_\beta^{\mathrm{inv}}=N\), so its contribution to \(\LX\) is
one.  The formula for \(\PX\) follows from
\eqref{eq:P}, and \cref{lem:comparison} bounds
\(N\rho_q(N)\) between \(1\) and \(3\).
\end{proof}

\begin{proposition}[Long fibers cannot return the endpoints]
\label{prop:endpoint-obstruction}
For every \(M\ge1\) there is an abstract nonnegative census with
unit atom cap and \(M\) one-point branches such that
\[
 \LX=0,\qquad\EX=\PX=M.
\]
\end{proposition}

\begin{proof}
Take \(M\) branches with \(N_\beta=1\), one atom of weight one on
each, and apply \cref{lem:endpoints}.
\end{proof}

\begin{remark}[Scope of the countermodels]
The two constructions are \(\Lzero\) sharpness examples for the
information used in a positive inequality.  They do not assert
that the actual TPC carrier contains such a family and are not stop
certificates for the physical route.  They prove that an actual
mass or branch-incidence theorem is still needed.
\end{remark}

\section{Finite exact certificate}
\label{sec:certificate}

The script
\texttt{experiments/tpc104\_principal\_mass\_audit.py} checks, using
exact rational arithmetic:
\begin{itemize}[leftmargin=2em]
\item every \(1\le N\le q\) for a deterministic list of odd
primes;
\item the endpoint identities and both inequalities in
\eqref{eq:comparison};
\item the exact disintegration on deterministic weighted branch
families; and
\item the two countermodel formulas.
\end{itemize}
Its JSON output is a finite regression certificate, not numerical
evidence that the actual growing ledgers satisfy their target.

\section{Route decision and claim boundary}
\label{sec:boundary}

\subsection*{Established}

\begin{itemize}[leftmargin=2em]
\item The \(N=1,2\) invertible mass is returned exactly as a purely
principal term.
\item The long principal coefficient is uniformly and sharply
comparable to \(1/N_\beta\).
\item Gate P is equivalent to the pair of actual weighted bounds
\(\EX,\LX\le X^{o(1)}Q_X^2\).
\item The TPC-103 cap reduces crude control of those ledgers to
actual branch censuses, but does not prove the censuses.
\end{itemize}

These are \(\Lzero\) coefficient identities and an \(\Lone\)
lossless reduction on the actual fixed-\(h_0\) carrier.

\subsection*{Not established}

\begin{itemize}[leftmargin=2em]
\item Neither \(\EX\) nor \(\LX\) is bounded at the target scale
for the growing physical packet.
\item No polynomial lower bound for either actual ledger is proved,
so the positive route is not declared stopped.
\item No cross-map estimate, signed affine M\"obius cancellation,
full-block return, or strict endpoint closure is proved.
\item No \(\Ltwo\) fixed-power saving, parity breakthrough,
Hardy--Littlewood asymptotic, prime-pair lower bound, or twin-prime
conclusion follows.
\end{itemize}

The next exact task is to quotient identical affine maps without
losing their input profile or source provenance.  Only after that
quotient is lossless should the genuinely distinct-map incidence
be attacked, as prescribed by TPC-MVP1
\citep{WangTPCMVP1}.

\begingroup
\footnotesize
\raggedright
\bibliographystyle{plainnat}
\bibliography{references}
\endgroup

\end{document}
